% Clase 2 --- Producto cartesiano como rectángulo, caso continuo (§5.4).
% Ancla: "AB. Es mi primer intervalo. Y yo voy a dibujar mi segundo intervalo
% de modo vertical. CD [...] ¿Qué forma va a tener el producto cartesiano?
% Será rectangular [...] Y cuál será el área de dicho rectángulo? b-a por d-c."
\begin{center}
\begin{tikzpicture}
  % Ejes
  \draw[figaxis] (-0.6,0) -- (5.6,0) node[below left, figlblaux] {};
  \draw[figaxis] (0,-0.6) -- (0,4.4);

  % Rectángulo A x B
  \fill[figblue!8] (1,1) rectangle (4,3.2);
  \draw[figaux] (1,1) rectangle (4,3.2);

  % Intervalo A = [a,b] sobre el eje x
  \draw[figintervalo] (1,0) -- (4,0);
  \node[figptocerrado] at (1,0) {};
  \node[figptocerrado] at (4,0) {};
  \node[figlbl, below] at (1,-0.05) {$a$};
  \node[figlbl, below] at (4,-0.05) {$b$};

  % Intervalo B = [c,d] sobre el eje y
  \draw[figintervalo] (0,1) -- (0,3.2);
  \node[figptocerrado] at (0,1) {};
  \node[figptocerrado] at (0,3.2) {};
  \node[figlbl, left] at (-0.05,1) {$c$};
  \node[figlbl, left] at (-0.05,3.2) {$d$};

  % Punto genérico (x0,y0) y guías punteadas
  \node[figptoY] (P) at (2.7,2.2) {};
  \node[figlbl, above right=1pt] at (P) {$(x_0,y_0)$};
  \draw[figaux] (P) -- (2.7,0) node[figptoY, below, inner sep=1.5pt] {};
  \draw[figaux] (P) -- (0,2.2) node[figptoY, left, inner sep=1.5pt] {};
  \node[figlblaux, below] at (2.7,-0.35) {$x_0$};
  \node[figlblaux, left] at (-0.35,2.2) {$y_0$};

  \node[figlbl] at (2.5,3.75) {$A\times B$};
\end{tikzpicture}\\[0.5em]
{\small\itshape Cuando $A=[a,b]$ y $B=[c,d]$ son intervalos, $A\times B$ es el
rectángulo cuyos lados son esos intervalos: el área $(b-a)(d-c)$ es, de nuevo,
un producto. Un punto genérico $(x_0,y_0)$ del rectángulo se proyecta (líneas
punteadas) sobre cada eje: su primera componente cae en $A$ y la segunda en
$B$.}
\end{center}
